% Сборка из корня репозитория: latexmk exam/exam-program.tex
\documentclass[11pt,a4paper]{article}
\newcommand{\examtitle}{Программа зачёта}
\usepackage{iftex}
\RequirePDFTeX
\usepackage{cmap}
\usepackage[T2A]{fontenc}
\usepackage[utf8]{inputenc}
\usepackage[english,russian]{babel}
\usepackage{PTSerif}
\usepackage{amsmath,amssymb,microtype}
\usepackage[a4paper,left=23mm,right=23mm,top=20mm,bottom=21mm]{geometry}
\usepackage{xcolor,enumitem,fancyhdr,hyperref}

\definecolor{accent}{HTML}{187AA5}
\definecolor{ink}{HTML}{302B28}
\definecolor{muted}{HTML}{6D6560}
\definecolor{tint}{HTML}{EDF4F6}
\definecolor{rulecolor}{HTML}{A9CADA}
\hypersetup{unicode=true,colorlinks=true,urlcolor=accent,linkcolor=accent,
  pdftitle={\examtitle},pdfauthor={Чиняев Владимир Николаевич}}
\color{ink}
\setlength{\parindent}{0pt}
\setlength{\parskip}{3pt}
\linespread{1.08}
\setlist[enumerate]{leftmargin=*,label=\textcolor{accent}{\arabic*.},
  widest=99,itemsep=3pt,topsep=4pt,parsep=0pt}
\setlength{\emergencystretch}{1em}
\pagestyle{fancy}
\fancyhf{}
\renewcommand{\headrulewidth}{0pt}
\renewcommand{\footrulewidth}{0.4pt}
\renewcommand{\footrule}{\hbox to\headwidth{\color{rulecolor}\leaders\hrule height\footrulewidth\hfill}}
\fancyfoot[L]{\footnotesize\color{muted}Введение в маневрирование КА\quad /\quad \examtitle}
\fancyfoot[R]{\footnotesize\color{muted}\thepage}

\newcommand{\examheading}[1]{%
  {\small\color{accent}ВВЕДЕНИЕ В МАНЕВРИРОВАНИЕ КА}\par
  \vspace{3pt}
  {\huge\examtitle}\par
  {\small\color{muted}#1}\par
  \vspace{4pt}
  {\color{rulecolor}\hrule}
  \vspace{5pt}}

\newcommand{\examnote}[1]{%
  \par\smallskip\noindent
  \begingroup
  \setlength{\fboxsep}{7pt}%
  \setlength{\fboxrule}{0.4pt}%
  \fcolorbox{rulecolor}{tint}{%
    \begin{minipage}{\dimexpr\linewidth-2\fboxsep-2\fboxrule\relax}
      \small #1\par
    \end{minipage}}%
  \endgroup\par\smallskip}

\newcommand{\topicheading}[1]{%
  \par\vspace{5pt}{\large\bfseries\color{accent}#1}\par\nobreak\smallskip}


\begin{document}
\examheading{Дифференцированный зачёт \enspace / \enspace 2026--2027 учебный год}

\examnote{\textbf{Правила проведения зачёта.} Студенту предлагаются три пункта
программы в случайном порядке. За первые два ответа можно получить до 2 баллов
за каждый, за третий~--- до 1 балла. Зачёт проходит без вспомогательных
материалов и устройств. В ответе достаточно объяснить смысл условий, моделей
и результатов, показать своё понимание.\par
\smallskip
\textbf{Итоговая оценка.} При сданном теоретическом минимуме суммируются баллы
за основные домашние задания (до 5), задачи со звёздочкой (до 1) и зачёт
(до 5). Итог ограничен 10 баллами. Если минимум не сдан, итоговая оценка~--- 2.}

\topicheading{Орбита и постановка задачи маневрирования}
\begin{enumerate}[series=topics]
\item \textbf{Классические кеплеровы элементы орбиты.} Радиусы апсид, эпоха
элементов. Классификация орбит по высоте и эксцентриситету. Условия геостационарности.
\item \textbf{Фаза и представление состояния.} Истинная, эксцентрическая и средняя
аномалии, смысл уравнения Кеплера. Вырождения классических элементов,
аргумент широты, декартовы элементы.
\item \textbf{Задача маневрирования.} Модели сил, КА и двигательной установки,
граничные условия, ограничения и критерий качества. Целевые условия задач
перехода, встречи и уклонения.
\end{enumerate}

\topicheading{Баллистический прогноз}
\begin{enumerate}[resume=topics]
\item \textbf{Возмущения и усреднение.} Вековые, долго- и короткопериодические
изменения. Оскулирующие и средние элементы, быстрые и медленные переменные,
смысл одно- и двукратного усреднения.
\item \textbf{Описание изменений элементов.} Малые параметры и порядок теории
возмущений, вариация параметров и условие оскуляции. Что описывают формы
Лагранжа и Гаусса?
\item \textbf{Действие возмущающих сил.} Вековые эффекты $J_2$ для $\Omega,\omega,M$,
солнечно-синхронная орбита и критическое наклонение. Влияние атмосферы,
третьих тел и солнечного давления на орбитальные элементы.
\item \textbf{Методы прогноза.} Численный, аналитический и полуаналитический
подходы и источники погрешности. Теории Козая и Брауэра, связь TLE с SGP4,
роль средних элементов и короткопериодических поправок в DSST.
\end{enumerate}

\topicheading{Одноимпульсные манёвры в центральном поле}
\begin{enumerate}[resume=topics]
\item \textbf{Импульсная модель.} Скачок скорости при непрерывном положении.
Действие радиальной, трансверсальной и нормальной компонент на энергию
и орбитальный момент. Связь компонент с орбитальным базисом.
\item \textbf{Переход между двумя заданными орбитами.} Как искать общие точки
на компланарных и некомпланарных орбитах? Как определить импульс в общей точке
и выбрать точку с меньшими затратами?
\end{enumerate}

\newpage
\topicheading{Околокруговое приближение}
\begin{enumerate}[resume=topics]
\item \textbf{Опорная круговая орбита и линеаризация.} Цилиндрические координаты,
вращение базиса и физические компоненты скорости. Какие отклонения считаются
малыми и где применимы линейные уравнения движения?
\item \textbf{Пассивное движение.} Связь изменений радиуса, продольной фазы
и нормального отклонения с большой полуосью, вектором эксцентриситета
и плоскостью орбиты. Условие отсутствия продольного дрейфа.
\item \textbf{Отклик на импульс.} Как скачок скорости влияет на последующее
свободное движение? Как в околокруговом приближении по эффектам отдельных
импульсов оценивают результат всей последовательности?
\item \textbf{Радиальный импульс.} Как меняются большая полуось и вектор
эксцентриситета? Как ведёт себя фаза на последующем пассивном участке?
Как устроено сближение радиальными импульсами?
\item \textbf{Трансверсальный импульс.} Изменение большой полуоси и периода,
форма новой орбиты и накопление фазового сдвига. Как расположение импульсов
в апсидах влияет на коррекцию большой полуоси и эксцентриситета?
\item \textbf{Нормальный импульс.} Поворот плоскости и траектория относительно
опорной круговой орбиты. Где приложить импульс для перехода в заданную плоскость?
\end{enumerate}

\topicheading{Оптимальность импульсных манёвров}
\begin{enumerate}[resume=topics]
\item \textbf{Задача оптимизации.} Выбор числа импульсов, их фаз, направлений
и величин. Критерий суммарной характеристической скорости. Связь верхней оценки
числа импульсов с числом независимых граничных условий.
\item \textbf{Базис-вектор Лоудена.} Смысл его связи с граничными условиями.
Как направление и модуль базис-вектора связаны с ненулевыми импульсами?
Что означает условие стационарности по фазе во внутренней точке интервала?
\item \textbf{Достаточные условия глобального минимума.} Как проверить
граничные условия, направление импульсов и ограничение на модуль базис-вектора?
Что означает выход выбранного годографа за единичную сферу?
\end{enumerate}

\topicheading{Оптимальный переход на околокруговую орбиту}
\begin{enumerate}[resume=topics]
\item \textbf{Задача перехода.} Какие исходные данные и конечные условия
задают? Как годограф базис-вектора помогает выбрать места приложения
и направления импульсов?
\item \textbf{Компланарный переход.} Касающиеся, пересекающиеся и непересекающиеся
орбиты. Возможное число, знаки и допустимые фазы трансверсальных импульсов.
Обобщённый переход Гомана.
\item \textbf{Некомпланарный переход.} Геометрия узлового решения и условия
его выбора. Смысл невырожденного и особого случаев, решения с нулевыми
радиальными компонентами и порядок выбора типа решения.
\end{enumerate}

\newpage
\topicheading{Встреча, задача Ламберта и относительное движение}
\begin{enumerate}[resume=topics]
\item \textbf{Линейная задача встречи в элементах.} Какие конечные условия
задают? Как импульс связан с изменением элементов и зачем
переносить эти отклонения к моменту встречи? Почему изменение большой полуоси
влияет на конечную фазу?
\item \textbf{Расчёт последовательности импульсов встречи.} Смысл матрицы
манёвра, разрешимость и обусловленность. Как выбор моментов и числа
импульсов связан с минимумом суммарной характеристической скорости?
\item \textbf{Задача Ламберта.} Постановка задачи. Семейство коник,
зависимость времени перелёта от параметра дуги и выбор ветви.
Как решение задачи используют для расчёта импульсов встречи?
\item \textbf{Модель Клохесси--Уилтшира.} Координаты во вращающемся базисе
цели и условия линеаризации. Плоские и нормальные колебания, продольный
дрейф и его связь с начальными условиями.
\item \textbf{Относительное движение на эллиптической орбите.} Модель
Чаунера--Хемпеля и назначение решения Яманаки--Анкерсена. Роль истинной аномалии
цели и связь с круговым случаем.
\item \textbf{Нелинейная проверка встречи.} Как по прогнозу оценить невязки
положения и скорости? Что меняют при ньютоновском уточнении импульсов
и как проверяют результат?
\end{enumerate}

\topicheading{Протяжённые манёвры и возмущения}
\begin{enumerate}[resume=topics]
\item \textbf{Модель конечной тяги.} Тяговое ускорение, расход массы,
характеристическая скорость и длительность включения. Как направление
тяги влияет на энергию и орбитальный момент?
\item \textbf{Замена импульса протяжённым включением.} Начало, центр и угловая
протяжённость манёвра. Смысл центрирования и влияние профиля ускорения.
Почему постоянное направление в инерциальной и орбитальной системах
даёт разные орбитальные результаты?
\item \textbf{Распределение манёвра между витками.} Смысл весов отдельных
включений, равные и убывающие веса. Масса, длительность и положение
манёвров. Последовательный прогноз активных и пассивных участков.
\item \textbf{Учёт возмущений и уточнение перехода.} Как $J_2$ меняет место
и направление следующих включений? Как сравнить достигнутую и целевую
орбиты в общей эпохе и использовать невязку для пристрелки параметров?
\end{enumerate}

\topicheading{Оценка опасного сближения}
\begin{enumerate}[resume=topics]
\item \textbf{Геометрия сближения.} Момент наибольшего сближения (TCA),
относительная скорость, плоскость сближения и вектор промаха. Что учитывают
суммарный радиус объектов и допущение о коротком пролёте?
\item \textbf{Вероятность столкновения.} Перенос ковариаций к моменту сближения,
условие их сложения и проекция в плоскость сближения. Зависимость вероятности
от промаха и эллипса ошибок. Размывание вероятности.
\end{enumerate}

\newpage
\topicheading{Манёвр уклонения}
\begin{enumerate}[resume=topics]
\item \textbf{Действие импульса на сближение.} Как манёвр меняет промах
и момент наибольшего сближения? Как его эффект зависит от направления
импульса и времени его приложения?
\item \textbf{Выбор манёвра уклонения.} Критерии сравнения вариантов и роль
расстояния Махаланобиса. Постановка задачи уклонения при свободном и заданном
направлении импульса.
\item \textbf{Нелинейная проверка манёвра.} Как по нелинейному прогнозу
проверяют рассчитанный манёвр? Какие характеристики сближения определяют
заново и по какому критерию оценивают результат?
\end{enumerate}

\topicheading{Выведение КА на орбиту}
\begin{enumerate}[resume=topics]
\item \textbf{Целевое состояние и программа выведения.} Условия выхода
на круговую орбиту. Как устроена программа выведения и для чего нужен
гравитационный разворот? Как в плоской модели задают направления скорости и тяги?
\item \textbf{Тяга, масса и ступенчатость.} Связь удельного импульса, расхода
топлива и характеристической скорости. Что происходит с массой и скоростью
при идеализированном отделении ступени и как это влияет на выведение?
\item \textbf{Скоростной баланс и ограничения.} Какие потери учитывают
в скоростном балансе? Как оценивают аэродинамические и тяговые нагрузки?
Как программа разворота влияет на потери и нагрузки?
\item \textbf{Пассивный подъём и довыведение.} Зачем нужен пассивный участок?
Как выбирают точку довыведения в импульсном приближении и момент начала
протяжённого включения? Как уточняют программу по результатам расчёта?
\item \textbf{Вращение Земли и геометрия запуска.} Как вращение Земли влияет
на инерциальную скорость ракеты и её скорость относительно воздуха?
Какие наклонения достижимы при прямом запуске? Как выбирают время старта
для выхода в заданную плоскость?
\end{enumerate}

\topicheading{Практические аспекты орбитальных операций}
\begin{enumerate}[resume=topics]
\item \textbf{Ошибки выведения и наведения.} Как ошибка выведения влияет
на дальнейший переход? Чем различаются ошибка наведения и ошибка оценки
ориентации? От чего зависят последствия ошибки наведения для орбиты?
\item \textbf{Отклонения и отказы двигательной установки.} Как уровень тяги
и удельный импульс влияют на длительность работы и расход рабочего тела?
Что меняют ошибки соотношения компонентов, разброс между включениями,
систематическое смещение и отказ двигателя?
\item \textbf{Нештатное выполнение программы.} Чем различаются последствия
опоздания, пропуска и аварийного прекращения манёвра? Что проверяют перед
выбором продолжения?
\item \textbf{Цикл орбитальных операций.} Каковы его основные стадии и для
чего нужна каждая? Как используют резервы? Какие задачи решают при
реконструкции выполненного манёвра и калибровке двигателя?
\end{enumerate}
\end{document}
